\clearpage\thispagestyle{empty}
\pdfbookmark[0]{Quellen zur Besprechung von Pasquichs Epitome}{pasquich1811-sources}
\begin{center}\textbf{Quellen zur Besprechung von Pasquichs Epitome}\end{center}
\noindent Diese KI-übersetzte Fassung in modernem Deutsch folgt Carl Friedrich Gauß, \textit{Werke}, Band~VI (Göttingen, 1874), S.~535--536. Die Besprechung erschien zuerst unter Wien in den \textit{Göttingischen gelehrten Anzeigen}, Stück~144, 9. September 1811, S.~1433--1436.
\par\smallskip\noindent Pasquichs \textit{Epitome elementorum astronomiae sphaerico-calculatoriae} erschien in zwei Teilen mit einem Anhang (Wien, 1811). Der Erstdruck der Besprechung und Pasquichs Titelblätter nennen den Verlag Schaumburg. Die Werkausgabe nennt beim Hauptwerk Schaumann und beim Anhang Schaumburg.
\par\smallskip\noindent Im Anhang gibt \href{https://api.digitale-sammlungen.de/iiif/image/v2/bsb10049414_00029/full/1250,/0/default.jpg}{Tafel XXVIII} eine Korrektur für Mittag und Mitternacht aus korrespondierenden Sonnenhöhen an. Gauß erwähnt in der Besprechung die Mittagskorrektur.
\par\smallskip\noindent \href{https://gdz.sub.uni-goettingen.de/id/PPN236007467}{Text der Werkausgabe}. Erstdruck, \href{https://images.p.sub.uni-goettingen.de/iiif/image/gdz:PPN319721507_1811:00001467/full/1650,/0/default.jpg}{S.~1433}, \href{https://images.p.sub.uni-goettingen.de/iiif/image/gdz:PPN319721507_1811:00001468/full/1650,/0/default.jpg}{S.~1434}, \href{https://images.p.sub.uni-goettingen.de/iiif/image/gdz:PPN319721507_1811:00001469/full/1650,/0/default.jpg}{S.~1435} und \href{https://images.p.sub.uni-goettingen.de/iiif/image/gdz:PPN319721507_1811:00001470/full/1650,/0/default.jpg}{S.~1436}. Titelblätter aus der Bayerischen Staatsbibliothek, \href{https://api.digitale-sammlungen.de/iiif/image/v2/bsb10049412_00005/full/1250,/0/default.jpg}{Teil I}, \href{https://api.digitale-sammlungen.de/iiif/image/v2/bsb10049413_00005/full/1250,/0/default.jpg}{Teil II} und \href{https://api.digitale-sammlungen.de/iiif/image/v2/bsb10049414_00003/full/1250,/0/default.jpg}{Anhang}.
\par\smallskip\noindent Übersetzung und Satz mit KI erstellt.
\clearpage
